% !PRINTED 671
\weberpage{671}
%<<CONT
in corresponding fashion take up the points at which $\eta$ has a negative order number. Then the polygons $\mathfrak{A}$ and $\mathfrak{B}$ are of equal order, namely of the order of the function $\eta$, and by these polygons the function $\eta$ is determined up to a constant factor.

We can therefore introduce the symbolic notation:

\weq{1}{\eta = \dfrac{\mathfrak{A}}{\mathfrak{B}}\,,}

and call $\mathfrak{A}$ the \gesperrt{numerator} or the \gesperrt{upper polygon}, $\mathfrak{B}$ the \gesperrt{denominator} or the \gesperrt{lower polygon} of $\eta$. According to this definition, $\mathfrak{A}$ and $\mathfrak{B}$ are in the first instance relatively prime to one another. But we wish to extend this notation still further by setting, when $\mathfrak{M}$ is an arbitrary polygon:

\weq{2}{\dfrac{\mathfrak{M}\mathfrak{A}}{\mathfrak{M}\mathfrak{B}} = \dfrac{\mathfrak{A}}{\mathfrak{B}}}

Then the notation (1) is free of the restriction that $\mathfrak{A}$ and $\mathfrak{B}$ are to be relatively prime. Both polygons are still of equal order, but this order can be greater than the order of $\eta$.

With this representation of the functions $\eta$ the same rules then hold for multiplication and division as in reckoning with numerical fractions in the region of the natural numbers.

In (1) $\mathfrak{A}, \mathfrak{B}$ cannot be arbitrary $m$-gons, and the investigation of the relation between them is the great question which, in another way, is answered by \gesperrt{Abel's theorem}.

We set up the following definition:

1. \gesperrt{If two $n$-gons $\mathfrak{A}, \mathfrak{A}'$ can be the upper polygon and the lower polygon of a function $\eta$ in $\Omega$, then $\mathfrak{A}$ and $\mathfrak{A}'$ are called equivalent.}

In symbols: if there is a function $\eta$ to which the notation

\[
\eta = \dfrac{\mathfrak{A}}{\mathfrak{A}'}
\]

belongs, then

\weq{3}{\mathfrak{A} \sim \mathfrak{A}'.}

From the formula

\[
\eta = \dfrac{\mathfrak{A}}{\mathfrak{A}'}\,, \quad \eta' = \dfrac{\mathfrak{A}}{\mathfrak{A}''}\,, \quad \dfrac{\eta'}{\eta} = \dfrac{\mathfrak{A}'}{\mathfrak{A}''}
\]

there results the theorem:
%CONT>>
